\originalsection{第2次习题课}
\sourceinfo{MIT 18.04, Spring 2018, Recitation 2; 题面 \nolinkurl{mit18_04s18_recit2-handout.pdf}, PDF第1页; 解答 \nolinkurl{mit18_04s18_recit2-solutions.pdf}, PDF第1页. 课程作者 Jeremy Orloff, 页内署名 Vishesh Jain, 2018, CC BY-NC-SA 4.0}
\begin{exercise}
\textbf{原题 Rec2.1。}\label{ass:rec2-1}
1.1. 直接用复导数定义证明 $f(z)=\bar z$ 不可复微分. 1.2. 当 $z\to0$ 时, 比值 $\bar z/z$ 可沿路径趋于哪些值?
\end{exercise}
\begin{solution}
据官方解答翻译补全. 在任一点 $z_0$, 差商为 $\overline h/h$. 沿 $h=t\in\R\setminus\{0\}$ 趋于0时等于1; 沿 $h=\ii t$ 时等于 $-1$. 因而导数在任何点都不存在. 对1.2, 写 $z=r\ee^{\ii\theta}$, 则 $\bar z/z=\ee^{-2\ii\theta}$; 沿固定辐角的射线趋于0可以取得单位圆上任一极限. 反过来, 比值的模恒为1, 所以任何存在的路径极限都在单位圆上. 这并不意味着二维极限存在.
\end{solution}
\begin{exercise}
\textbf{原题 Rec2.2。}\label{ass:rec2-2}
2.1. 求 $\cos z$ 与 $\sin z$ 的实部、虚部. 2.2. 当 $z=x+\ii0$ 及 $z=0+\ii y$ 时分别是什么? 2.3. 它们在复平面上有界吗? 2.4. $\cos^2z+\sin^2z=1$ 是否成立?
\end{exercise}
\begin{solution}
据官方解答翻译补全. 由指数定义得到
\[
\cos(x+\ii y)=\cos x\cosh y-\ii\sin x\sinh y,\qquad
\sin(x+\ii y)=\sin x\cosh y+\ii\cos x\sinh y.
\]
故实轴上分别为 $\cos x$, $\sin x$, 虚部均为0; 虚轴上分别为 $\cosh y$, $\ii\sinh y$. 当 $y\to\infty$ 时 $\cosh y$ 与 $\sinh y$ 无界, 两个函数在 $\C$ 上均无界. 最后展开两个平方, 虚部相消, 实部为 $(\cos^2x+\sin^2x)(\cosh^2y-\sinh^2y)=1$, 故恒等式仍成立.
\end{solution}
\begin{exercise}
\textbf{原题 Rec2.3。}\label{ass:rec2-3}
3.1. 证明 $\ee^z$ 是 $z$ 的连续函数. 3.2. 据此证明 $\cos z$ 与 $\sin z$ 连续. 3.3. $\bar z$ 连续吗? 是否与第1题一致?
\end{exercise}
\begin{solution}
据官方解答翻译补全. $\ee^{x+\ii y}=\ee^x\cos y+\ii\ee^x\sin y$ 的两个实坐标均是连续实函数的乘积, 因而指数函数连续. 用 $\cos z=(\ee^{\ii z}+\ee^{-\ii z})/2$ 和 $\sin z=(\ee^{\ii z}-\ee^{-\ii z})/(2\ii)$, 连续函数的复合与线性组合仍连续. 对共轭, $|\bar z-\bar z_0|=|z-z_0|$, 直接得到连续性. 复可微蕴含连续, 反向不成立, 所以它处处连续与处处不可复微分没有矛盾.
\end{solution}
\begin{exercise}
\textbf{原题 Rec2.4。}\label{ass:rec2-4}
4.1. 把 $\ee^z,z^2,\cos z,\sin z$ 写成 $u(x,y)+\ii v(x,y)$. 4.2. 计算每个函数的 $u_x,u_y,v_x,v_y$, 检验 Cauchy--Riemann 方程. 4.3 (原文第一次编号). 求各函数的 $f'(z)$. 4.3 (原文第二次编号). 对 $\bar z$ 重复以上步骤, 与第1题是否一致? 为避免原文重复编号产生歧义, 两问分别记为 4.3-first 与 4.3-second.
\end{exercise}
\begin{solution}
编者补充 (AI): 官方解答仅要求参照讲义, 此处给出全部计算. \textbf{4.1--4.2.} 对 $\ee^z$, $u=\ee^x\cos y$, $v=\ee^x\sin y$, 从而 $(u_x,u_y,v_x,v_y)=(\ee^x\cos y,-\ee^x\sin y,\ee^x\sin y,\ee^x\cos y)$. 对 $z^2$, $u=x^2-y^2$, $v=2xy$, 四个偏导为 $(2x,-2y,2y,2x)$. 对 $\cos z$, $u=\cos x\cosh y$, $v=-\sin x\sinh y$, 四个偏导为
\[
(-\sin x\cosh y,\ \cos x\sinh y,\ -\cos x\sinh y,\ -\sin x\cosh y).
\]
对 $\sin z$, $u=\sin x\cosh y$, $v=\cos x\sinh y$, 四个偏导为 $(\cos x\cosh y,\sin x\sinh y,-\sin x\sinh y,\cos x\cosh y)$. 它们均处处满足 $u_x=v_y$, $u_y=-v_x$, 且偏导连续, 因而均为整函数.

\textbf{4.3-first.}\label{ass:rec2-4-3-first} 用 $f'=u_x+\ii v_x$ 依次得到 $\ee^z,2z,-\sin z,\cos z$. \textbf{4.3-second.}\label{ass:rec2-4-3-second} 对共轭函数 $u=x$, $v=-y$, 四个偏导为 $(1,0,0,-1)$; 第一条 Cauchy--Riemann 方程不成立, 所以无一处可复微分, 与第1题一致.
\end{solution}
